% 原书分章 PDF 第 12—13 页：第 25.3.3 节。
\subsection{作为 \texorpdfstring{$\psi_0$}{ψ0} 估计量的 CV-TMLE 的渐近正态性}

令 $P_0\in\mathcal M(J_0)$。
下面的定理直接来自
$(nh^d)^{0.5}(\psi_n^*-\psi_{h0})\Rightarrow_dN(0,\sigma_0^2)$
与 $h^{-J_0}(\psi_{h0}-\psi_0)\to B_0(J_0)$，二者表明
\[
 \psi_n^*-\psi_0\approx
 N(h^{J_0})B_0(J_0),\sigma_0^2/(nh^d)).
\]
平衡偏差平方与方差，得到最优速率
$h_n^*=n^{-1/(2J_0+d)}$；略微欠平滑（undersmoothing），则得到均值为零的极限分布 $N(0,\sigma_0^2)$。

\begin{bookresult}{定理}{25.2}
假定采用 $J$ 阶正交核 $K$，并满足定理 25.1 的条件。
令 $P_0\in\mathcal M(J_0)$，其中光滑程度 $J_0\geq1$ 未知。
令 $J_0^*=\min(J_0,J)$。
带宽的最优速率为 $h_n^*=n^{-1/(2J_0^*+d)}$。
若选择这一速率，则得到
\[
 n^{-J_0^*/(2J_0^*+d)}(\psi_n^*-\psi_0)
 \Rightarrow_dN(B_0(J_0^*),\sigma_0^2).
\]
若选择 $h_n$ 使 $h_n/n^{-1/(2J_0^*+d)}\to0$，即进行欠平滑，则得到
\[
 \sqrt{nh_n^d}(\psi_n^*-\psi_0)\Rightarrow_dN(0,\sigma_0^2).
\]
若选择这样的欠平滑带宽 $h_n$，则 $\psi_0$ 的渐近有效的 0.95 置信区间为
\[
 \psi_n^*\pm1.96\sigma_n/(nh_n^d)^{0.5},
\]
其中 $\sigma_n^2$ 为 $\sigma_0^2$ 的相合估计量。
\end{bookresult}

若已知 $J_0\geq J$，则上述最优速率 $h_n^*$ 已知，由 $J_0^*=J$ 确定。
但在不知道这一点的情况下，最优速率 $h_n^*$ 未知。
即使 $J_0\geq J$，从有限样本表现来看，我们仍希望选择
$ch_n^*$，其中 $c$ 应取一个合适的值，而这一值未知。
因此，无论哪一种情形，提出一个数据自适应的 $h$ 选择器都很重要。下一节将给出这样的选择器。
